\section{Radial Parts and the Moyal Normalization}
\label{sec:app-radial-parts-moyal}

This appendix records the proof spine for
Theorem~\ref{thm:lane-moyal-degree-zero}.  The stable core is the
scalar-reduced Weyl/Moyal normalization, the Capelli change of basis,
and the stable connected trace bracket.  The finite-\(N\) all-\(N\)
radial-parts step is conditional on the
Harish-Chandra--Wallach--Levasseur--Stafford radial-parts input in the
normalization stated below.  The Costello brane-defect graph
realization is not part of this appendix.

\begin{defn}[Formal Weyl/Moyal convention]
\label{def:app-radial-moyal-convention}
  Let
  \[
    P=\partial_{z_1}\otimes\partial_{z_2}
      -\partial_{z_2}\otimes\partial_{z_1}.
  \]
  The formal Weyl product on \(\C[[z_1,z_2]][[\hbar]]\) is
  \[
    f*g=m\circ\exp\left(\frac{\hbar}{2}P\right)(f\otimes g),
  \]
  and
  \[
    [f,g]_{\mathrm{raw}}=f*g-g*f,\qquad
    \{f,g\}_\hbar=\hbar^{-1}[f,g]_{\mathrm{raw}}.
  \]
  The reduced quantum Hamiltonian Lie algebra is the scalar quotient
  \[
    \bar A_\hbar=\C[[z_1,z_2]][[\hbar]]/\C[[\hbar]]
  \]
  as a Lie algebra.  It is not used as an associative quotient.
\end{defn}

\begin{prop}[Moyal coefficients]
\label{prop:app-radial-moyal-coefficients}
  For \(f=z_1^kz_2^l\) and \(g=z_1^mz_2^n\),
  \[
    P^r(f,g)=
    \sum_{a=0}^r(-1)^a\binom ra
      (k)_{r-a}(l)_a(m)_a(n)_{r-a}
      z_1^{k+m-r}z_2^{l+n-r}.
  \]
  Consequently
  \[
    [f,g]_{\mathrm{raw}}
      =\sum_{\substack{r\geq1\\r\ \mathrm{odd}}}
        \frac{2}{2^r r!}\hbar^rP^r(f,g),
  \]
  so the first and third raw commutator coefficients are
  \[
    \hbar P(f,g),\qquad \frac{\hbar^3}{24}P^3(f,g),
  \]
  and all even raw commutator coefficients vanish.
\end{prop}

\begin{proof}
  The two summands in
  \(P=\partial_{z_1}\otimes\partial_{z_2}
     -\partial_{z_2}\otimes\partial_{z_1}\)
  commute as differential operators on the tensor product.  The
  binomial theorem gives
  \[
    P^r=\sum_{a=0}^r(-1)^a\binom ra
      \partial_{z_1}^{r-a}\partial_{z_2}^{a}
      \otimes
      \partial_{z_1}^{a}\partial_{z_2}^{r-a}.
  \]
  Applying this to the two monomials gives the falling-factorial
  formula.  Interchanging \(f\) and \(g\) changes \(P^r(f,g)\) by
  \((-1)^r\), so the commutator retains exactly the odd powers and
  contributes the coefficient \(2/(2^r r!)\) in order \(r\).
\end{proof}

\begin{defn}[Matrix Weyl algebra and Weyl traces]
\label{def:app-radial-matrix-weyl}
  Let \(\mathcal W_N\) be the completed Weyl algebra generated by
  entries of two \(N\times N\) matrices \(X,Y\) with
  \[
    [X^i{}_j,Y^k{}_l]=\hbar\,\delta^i_l\delta^k_j,\qquad
    [X^i{}_j,X^k{}_l]=[Y^i{}_j,Y^k{}_l]=0.
  \]
  The normal-ordered quantum comoment for conjugation is
  \[
    \widehat\mu_\epsilon=-\operatorname{Tr}([\epsilon,X]Y),
    \qquad
    \hbar^{-1}[\widehat\mu_\epsilon,a]=\epsilon\cdot a.
  \]
  For \(f\in\C[z_1,z_2][[\hbar]]\), define the Weyl-ordered trace
  \[
    J_N(f)=\operatorname{Tr}\operatorname{Op}_{\mathrm W}(f)(X,Y).
  \]
  The scalar line is removed after passing to the degree-zero quantum
  Hamiltonian reduction
  \[
    \mathrm N_{\mathcal W_N}
      (\mathcal W_N\widehat\mu(\mathfrak{gl}_N))/
    \mathcal W_N\widehat\mu(\mathfrak{gl}_N).
  \]
\end{defn}

\begin{lem}[Capelli--Weyl triangular normalization]
\label{lem:app-radial-capelli-triangular}
  Let \(T_{a,b}=\operatorname{Tr}(X^aY^b)\) be the normal-ordered trace.
  In the quantum Hamiltonian reduction, the locally proved
  Weyl-normal-ordering formula is
  \[
    J_N(z_1^az_2^b)\equiv
      \sum_{r=0}^{\min(a,b)}
        \left(-\frac{\hbar N}{2}\right)^r
        r!\binom ar\binom br\,T_{a-r,b-r}.
  \]
  The inverse triangular system is
  \[
    T_{a,b}\equiv
      \sum_{r=0}^{\min(a,b)}
        \left(\frac{\hbar N}{2}\right)^r
        r!\binom ar\binom br\,J_N(z_1^{a-r}z_2^{b-r}).
  \]
\end{lem}

\begin{proof}
  Total Weyl symmetrisation is obtained from normal ordering by
  commuting \(Y\)'s past \(X\)'s.  Each contraction contributes the
  Capelli contact term coming from
  \[
    YX-XY+\hbar N I\equiv0
      \pmod{\mathcal W_N\widehat\mu(\mathfrak{gl}_N)}.
  \]
  Choosing \(r\) contracted \(X\)-slots and \(r\) contracted \(Y\)-slots
  gives \(\binom ar\binom br\), pairing them gives \(r!\), and Weyl
  ordering contributes \(2^{-r}\).  The sign is negative because the
  displayed moment relation rewrites a normal-ordering swap as the
  lower triangular term \(-\hbar N\).  The diagonal term is \(r=0\),
  hence the system is unipotent over \(\C[[\hbar N]]\); inversion
  changes \(-\hbar N/2\) to \(+\hbar N/2\).

  The cited Capelli literature supplies the standard source for the
  Capelli contact phenomenon.  No primary source is used here for this
  exact bivariate trace formula; the displayed contraction count is the
  proof of the normalization in this manuscript.
\end{proof}

\begin{stmt}[External radial-parts input]
\label{stmt:app-radial-external-input}
  The finite-\(N\) radial-parts step uses the following external
  Harish-Chandra radial-parts theorem for invariant differential
  operators, in the form refined by Wallach and by
  Levasseur--Stafford
  \cites{harish-chandra-invariant-differential-operators,
  wallach-reductive-invariant-differential,
  levasseur-stafford-harish-chandra,
  levasseur-stafford-kernel-harish-chandra}.  After identifying
  \(\mathcal W_N\simeq\mathcal D_\hbar(\mathfrak{gl}_N)\), the
  Harish-Chandra map on invariants is assumed to have the quantum
  moment-map ideal as kernel, hence to inject the quantum Hamiltonian
  reduction by conjugation into Weyl-invariant differential operators on
  the regular Cartan.  With
  \(\lambda_1,\ldots,\lambda_N\) the Cartan coordinates and
  \(\Delta=\prod_{i<j}(\lambda_i-\lambda_j)\), the radial-part map is
  assumed to be conjugated by \(\Delta\).  This appendix treats this
  precise Vandermonde-conjugated compatibility statement as source
  input; it is not re-proved here.

  The input includes the exact image normalization for the Weyl trace
  generators used in this paper:
  \[
    \operatorname{Rad}_0(J_N(f))
      =\Delta^{-1}\left(\sum_{i=1}^N
        \operatorname{Op}_{\mathrm W}
        (f)(\lambda_i,-\hbar\partial_{\lambda_i})\right)\Delta.
  \]
  Lemma~\ref{lem:app-radial-pure-shear-extension} proves that a smaller
  source package would suffice: the pure \(z_1\)- and \(z_2\)-power
  images together with compatibility of the radial map with the Weyl
  adjoint shears.  The
  Harish-Chandra--Wallach--Levasseur--Stafford anchors presently verify
  the homomorphism, regular-Cartan landing, and kernel/surjectivity
  context; they do not by themselves verify this exact Weyl-trace image
  convention.  That convention remains the named conditional input.
\end{stmt}

\begin{lem}[Pure powers and shears determine the image normalization]
\label{lem:app-radial-pure-shear-extension}
  Fix \(N\).  Suppose a Vandermonde-conjugated radial map is defined on
  the reduced invariant quotient and is compatible with the following
  data:
  \[
    \operatorname{Rad}_0(J_N(z_1^a))
      =\Delta^{-1}\Bigl(\sum_i\lambda_i^a\Bigr)\Delta,\qquad
    \operatorname{Rad}_0(J_N(z_2^b))
      =\Delta^{-1}\Bigl(\sum_i(-\hbar\partial_{\lambda_i})^b\Bigr)\Delta,
  \]
  for all \(a,b\geq0\), and for every polynomial \(p\) it intertwines
  the Weyl adjoint shear
  \[
    X\mapsto X,\qquad Y\mapsto Y+p(X)
  \]
  with the Cartan shear
  \[
    \lambda_i\mapsto\lambda_i,\qquad
    -\hbar\partial_{\lambda_i}\mapsto
      -\hbar\partial_{\lambda_i}+p(\lambda_i).
  \]
  With the sign convention \(Y=-\hbar\partial_X^t\), this is the plus
  shear generated by \(\operatorname{Tr}(X^{a+1})/(a+1)\).  Then the
  full image formula
  \[
    \operatorname{Rad}_0(J_N(f))=
      \Delta^{-1}\left(\sum_i
      \operatorname{Op}_{\mathrm W}
      (f)(\lambda_i,-\hbar\partial_{\lambda_i})\right)\Delta
  \]
  holds for every polynomial \(f\in\C[z_1,z_2]\).
\end{lem}

\begin{proof}
  It is enough to treat monomials.  The case \(z_1^a\) is one of the
  assumed pure-power images.  For \(b\geq0\) and \(a\geq0\), apply the
  shear with \(p_t(X)=tX^a\) to the pure generator \(J_N(z_2^{b+1})\).
  Weyl quantization is linear in the commutative symbol and the shear
  sends the symbol \(z_2\) to \(z_2+t z_1^a\), hence the coefficient of
  \(t\) in the matrix Weyl trace is
  \[
    (b+1)J_N(z_1^az_2^b).
  \]
  On the Cartan side the same shear sends
  \(-\hbar\partial_{\lambda_i}\) to
  \(-\hbar\partial_{\lambda_i}+t\lambda_i^a\), so the coefficient of
  \(t\) in
  \[
    \sum_i
    \operatorname{Op}_{\mathrm W}
    \bigl((z_2+t z_1^a)^{b+1}\bigr)
      (\lambda_i,-\hbar\partial_{\lambda_i})
  \]
  is
  \[
    (b+1)\sum_i
    \operatorname{Op}_{\mathrm W}
      (z_1^az_2^b)(\lambda_i,-\hbar\partial_{\lambda_i}).
  \]
  Intertwining of the shears and the pure \(z_2^{b+1}\)-image therefore
  gives the required image for \(z_1^az_2^b\), after division by
  \(b+1\).  Linearity gives the result for all polynomials.
\end{proof}

\begin{thm}[Finite-\(N\) radial-parts identity, conditional on the input]
\label{thm:app-radial-finite-N}
  Under the radial-parts input of
  Statement~\ref{stmt:app-radial-external-input}, put
  \[
    S_N(f)=\sum_{i=1}^N
      \operatorname{Op}_{\mathrm W}
      (f)(\lambda_i,-\hbar\partial_{\lambda_i}).
  \]
  Then
  \[
    \operatorname{Rad}_0(J_N(f))=\Delta^{-1}S_N(f)\Delta,
  \]
  and for all \(f,g\)
  \[
    D_N(f,g)=
      \hbar^{-1}[J_N(f),J_N(g)]-J_N(\{f,g\}_\hbar)
  \]
  lies in the quantum moment-map ideal.  Thus \(D_N(f,g)\) vanishes in
  degree-zero quantum Hamiltonian reduction.
\end{thm}

\begin{proof}
  The first displayed identity is the exact image normalization included
  in Statement~\ref{stmt:app-radial-external-input}.  The assumed
  radial-parts input also makes the conjugated radial map injective on
  the reduced invariant quotient, so the commutator defect can be tested
  on the Cartan.

  Lemma~\ref{lem:app-radial-pure-shear-extension} explains the minimal
  normalization package that would imply the full displayed formula:
  pure \(z_1\)-powers, pure \(z_2\)-powers, and compatibility with the
  plus shear \(Y\mapsto Y+p(X)\) in the convention
  \(Y=-\hbar\partial_X^t\), equivalently
  \(z_2=-\hbar\partial_\lambda\) on the Cartan.  The present sources do
  not verify that package; it remains part of the external input.

  Finally the \(i\)-th summand of \(S_N\) is an ordinary one-particle
  Weyl algebra and distinct summands commute, hence
  \[
    \hbar^{-1}[S_N(f),S_N(g)]=S_N(\{f,g\}_\hbar).
  \]
  Conjugation by \(\Delta\) preserves commutators, so
  \(\operatorname{Rad}_0(D_N(f,g))=0\).  Injectivity of the radial
  map on the reduced quotient puts \(D_N(f,g)\) in the quantum
  moment-map ideal.
\end{proof}

\begin{cor}[Conditional stable scalar-reduced connected bracket]
\label{cor:app-radial-stable-connected}
  Under the same radial-parts input as
  Theorem~\ref{thm:app-radial-finite-N}, for each fixed polynomial
  degree the classes
  \[
    \overline{J_N(f)}
  \]
  have a stable connected large-\(N\) limit after removing the empty
  trace and disconnected products.  In that limit,
  \[
    \hbar^{-1}[\partial_{\mathrm{bb},\hbar}(f),
      \partial_{\mathrm{bb},\hbar}(g)]_{\mathrm{conn,raw}}
      =
      \partial_{\mathrm{bb},\hbar}(\{f,g\}_\hbar).
  \]
\end{cor}

\begin{proof}
  The Capelli system of Lemma~\ref{lem:app-radial-capelli-triangular}
  is triangular with leading term \(T_{a,b}\).  Its lower terms have
  smaller total degree and include the scalar empty trace only when
  both indices have been contracted away.  The connected projection
  removes the empty trace and disconnected products, while preserving
  the single-trace \(J_N\)-basis.  The finite-\(N\) identity of
  Theorem~\ref{thm:app-radial-finite-N} is compatible with adjoining
  extra Cartan variables, because \(S_{N+k}(f)\) is obtained from
  \(S_N(f)\) by adding commuting one-particle summands.  Therefore the
  bracket identity stabilizes degreewise in \(N\).
\end{proof}

\begin{rmk}[Scalar quotient and the \(U(1)\) anomaly]
\label{rmk:app-radial-scalar-quotient}
  Constants are central in the Moyal Lie algebra and are killed in
  \(\bar A_\hbar\).  Before this scalar reduction the linear
  commutator carries the central term
  \([z_1,z_2]_*= \hbar\), and on the brane side the corresponding
  matrix trace is the Capelli shift \(\hbar N\).  The conditional
  reduced theorem is the theorem after quotienting by this scalar line.
  The unreduced line is the \(U(1)\) anomaly class of
  Theorem~\ref{thm:u1-center-anomaly}.
\end{rmk}

\begin{rmk}[Relation to the finite exact verification]
\label{rmk:app-radial-finite-verification}
  The finite exact verification is independent evidence for the local
  Moyal and Capelli normalizations, not a substitute for the external
  radial-parts input. It uses rational arithmetic to check:
  \[
  \begin{gathered}
    \text{Moyal monomial pairs with exponents }0,\ldots,10
    \text{ through order }11,\\
    \text{the Capelli triangular round trip for }a,b\leq10,\\
    \text{the direct }N=2\text{ and }N=3\text{ radial restrictions on }
    80\text{ operator/test-polynomial pairs in each rank},\\
    [S_2(f),S_2(g)]=S_2([f,g]_*)
    \text{ on the primary and higher-order radial test pairs},\\
    \text{the leading connected-cumulant bracket in the primary cases},\\
    \text{and the open-line first/third graph weights.}
  \end{gathered}
  \]
  The finite exact verification therefore checks the coefficients \(1\) and \(1/24\), the
  vanishing of even commutator orders in the tested range, the
  Capelli signs, the rank-two Cartan shadow of the radial-parts
  identity, and the direct rank-two and rank-three restriction of
  \(J_N(f)\) on the stated \(80\) operator/test-polynomial pairs in
  each of ranks \(N=2\) and \(N=3\).  It does not prove
  injectivity of the all-\(N\) Harish-Chandra radial-part map or source
  the exact image normalization beyond these low-rank tests.
\end{rmk}

\begin{rmk}[Residual graph problem]
\label{rmk:app-radial-residual-graph-problem}
  The formal Weyl/Moyal product fixes the algebraic target.  It does
  not prove that a Costello closed-open brane-defect model realizes
  this target to all orders.  That requires the sector specialization,
  propagator, counterterm, anomaly, and QME data named in
  Problem~\ref{prob:first-third-graph-normalizations}.
\end{rmk}
